% ============================================================
% Section 5: Discussion
% BearMamba-3: Mamba-3 + L_kin + Multi-Sensor Fusion
% Target: MSSP (IF 8.4)  Draft: 2026-06-15
% ============================================================

\section{Discussion}
\label{sec:discussion}

\paragraph{Why BearMamba-3 in lieu of BearMamba-2.}
BearMamba-2 outperforms BearMamba-3 on CWRU under heavy AWGN
($\leq -4$~dB, $p = 0.031$, Table~\ref{tab:snr_ablation}).
This is attributed to BearMamba-3's more complex optimisation landscape
arising from complex-state parameters (D18): in simple i.i.d.\ Gaussian
noise at constant speed, BearMamba-3 converges more slowly and with higher
inter-seed variance.
On PU cross-condition data, both architectures reach equivalent accuracy
($p > 0.05$, Table~\ref{tab:c1_pu}), and BearMamba-3's variance is
$3\times$ higher than BearMamba-2 (0.72\% vs 0.24\%), consistent with
the same convergence instability.
We note this as a known cost of BearMamba-3 and motivate its adoption
exclusively on the structural binding to $\mathcal{L}_{\text{kin}}$:
BearMamba-2 lacks instantaneous rotation frequencies and therefore cannot
carry the physical frequency alignment constraint (Section~\ref{subsec:lkin}).

\paragraph{Pre-classifier feature space under heavy noise.}
Figure~\ref{fig:tsne_snr6} shows t-SNE of pre-classifier features at SNR~$=-6$~dB
for BM3 $\pm\mathcal{L}_{\text{kin}}$ (seed~2, representative of the 5-seed ensemble).
At this SNR, both conditions produce heavily overlapping clusters
(Inner, Ball, and Outer interleaved), and the DB index is slightly higher
(worse) with $\mathcal{L}_{\text{kin}}$ (0.909 vs 0.883).
This is consistent with the non-significant accuracy difference at SNR~$=-6$~dB
($\Delta_{\text{mean}} = -0.09$~pp, $p \geq 0.25$): $\mathcal{L}_{\text{kin}}$
organises the state frequency space (Figure~\ref{fig:tsne_fbar}) but cannot
recover the classifier-level separability that is lost to signal corruption.
The variance-stabilisation effect (std $\downarrow 14\%$ at SNR~$=-6$~dB,
Table~\ref{tab:snr_ablation}) reflects a more consistent optimisation trajectory,
not improved feature geometry.

\begin{figure}[htbp]
\centering
\includegraphics[width=0.92\linewidth]{tsne_cwru_snr-6_seed2}
\caption{t-SNE of pre-classifier features, CWRU SNR~$=-6$~dB, seed~2.
  \textbf{(a)} Without $\mathcal{L}_{\text{kin}}$: DB~$= 0.883$;
  \textbf{(b)} With $\mathcal{L}_{\text{kin}}$: DB~$= 0.909$.
  At heavy noise, fault-class clusters overlap substantially for both conditions;
  $\mathcal{L}_{\text{kin}}$ does not recover separability.
  The slightly higher DB for $+\mathcal{L}_{\text{kin}}$ is within the
  inter-seed variation range and does not indicate structural degradation.}
\label{fig:tsne_snr6}
\end{figure}

\paragraph{Cross-seed t-SNE robustness.}
To validate that the single-seed visualisation in Figure~\ref{fig:tsne_snr6}
is representative of the 5-seed ensemble rather than an isolated artefact,
we re-extracted pre-classifier features for all five seeds at both
SNR~$=0$~dB and SNR~$=-6$~dB, computed t-SNE projections under matched
hyperparameters, and report the Davies--Bouldin (DB) index per seed
in Figure~\ref{fig:tsne_multiseed}.
At SNR~$=0$~dB, DB~$= 0.396 \pm 0.031$ across seeds (range
0.346--0.423), confirming highly consistent feature-space organisation
with $\mathcal{L}_{\text{kin}}$ when the signal-to-noise ratio is moderate.
At SNR~$=-6$~dB, DB~$= 0.920 \pm 0.103$ (range 0.833--1.097); the
elevated dispersion is concentrated in a single outlier (seed~1,
DB~$= 1.097$, val.\ accuracy $89.21\%$), reflecting a poor random
initialisation that converged to a degraded local minimum under heavy
noise---a known phenomenon for complex-state SSMs (D18 in our
project log).
Excluding this outlier yields DB~$= 0.875 \pm 0.030$, comparable to the
SNR~$=0$~dB stability, while including it provides an honest assessment
of seed sensitivity at extreme noise.
Notably, seed~2 (used in Figure~\ref{fig:tsne_snr6}) reports
DB~$= 0.906$, within $1.4\%$ of the 5-seed mean, confirming its
representativeness.
This cross-seed quantitative evidence complements the per-seed
qualitative t-SNE in Figure~\ref{fig:tsne_snr6} and pre-empts a
potential reviewer concern that single-seed visualisations may
cherry-pick favourable initialisations.

\begin{figure*}[htbp]
\centering
\includegraphics[width=0.95\linewidth]{tsne_multiseed_composite}
\caption{Cross-seed t-SNE robustness validation.
  Pre-classifier feature projections for all five training seeds
  at SNR~$=0$~dB (top row) and SNR~$=-6$~dB (bottom row), with
  Davies--Bouldin (DB) index annotated per seed.
  At SNR~$=0$~dB, DB~$= 0.396 \pm 0.031$ (5 seeds), demonstrating
  high feature-space consistency.
  At SNR~$=-6$~dB, DB~$= 0.920 \pm 0.103$; dispersion is dominated
  by seed~1 (DB~$= 1.097$, val.\ accuracy $89.21\%$), a single poor
  initialisation. Excluding seed~1 yields DB~$= 0.875 \pm 0.030$,
  matching the SNR~$=0$~dB stability. Seed~2 (used in
  Figure~\ref{fig:tsne_snr6}) reports DB~$= 0.906$, validating its
  representativeness of the 5-seed mean ($1.4\%$ relative deviation).
  All seeds use identical t-SNE hyperparameters
  (perplexity $= 30$, random\_state $= 42$).}
\label{fig:tsne_multiseed}
\end{figure*}

\paragraph{$\mathcal{L}_{\text{kin}}$ boundary conditions.}
$\mathcal{L}_{\text{kin}}$ shows variance-stabilising effects on CWRU
(AWGN, variable SNR) but is ineffective on the Paderborn dataset (discrete
speed switching) and XJTU-SY (run-to-failure, amplitude-driven regime).
The Paderborn case degrades the per-sample RPM advantage to a two-point
prior: with only two nominal speeds (900 and 1500~rpm), $\mathcal{F}_i$
reduces to two fixed frequency sets and the continuous-speed adaptation
advantage is not activated~\cite{an2019rnn_bearing,liao2020variable_domain}.
The XJTU case saturates the coverage-distance metric from epoch~1 onward
(BPFO distance $< 5$~Hz before $\mathcal{L}_{\text{kin}}$ even converges),
because large-amplitude end-of-life impulsive features naturally span the
fault frequency band; the remaining gradient signal from $\mathcal{L}_{\text{kin}}$
is too small to influence accuracy-relevant features.
Mechanism-driven approaches such as Jiang et al.~\cite{jiang2025mechanism}
similarly note that the effectiveness of physical priors is bounded by
the signal characteristics — a consistent observation across
physics-informed diagnostic methods.
Future work could investigate task-adaptive gating of $\mathcal{L}_{\text{kin}}$
based on detected signal impulsiveness (e.g.\ kurtosis threshold), to
gracefully deactivate the constraint in amplitude-dominated regimes.

\paragraph{LOBO OR degradation in dual-sensor XJTU.}
The $-15.5$~pp LOBO OR macro-recall drop (Table~\ref{tab:xjtu}) warrants
careful interpretation.
Condition~3 has only two OR bearings (Bearing3\_1 and Bearing3\_5); each
LOBO fold therefore trains on a single OR instance.
The vertical acceleration component of a single OR bearing captures
idiosyncratic mounting and preload characteristics that do not generalise
across bearings~\cite{xjtu_paper}.
In contrast, horizontal components are more structurally correlated with
OR fault patterns, explaining the significantly smaller degradation in IR
folds where two training IR bearings are available.
This degradation is a dataset-structure constraint, not an inherent
limitation of dual-sensor architectures; a balanced dataset with
$\geq 4$ OR bearings per condition would be expected to eliminate this
effect.
Practitioners deploying multi-sensor fusion in LOBO-style evaluations
should ensure adequate per-fault-type bearing coverage in the training set.

\paragraph{Two complementary inductive bias paradigms.}
The cross-dataset analysis in Section~\ref{subsec:paradigm} exposes a
structural complementarity in how two families of methods achieve
low-SNR robustness.

The \textbf{BN-dependent paradigm} (1D-CNN, most CNN architectures) relies on
BatchNorm layers that normalise activations using source-domain batch statistics.
This provides powerful implicit regularisation within a fixed distribution:
1D-CNN reaches $95.96\%$ at CWRU $-8$~dB, outperforming BM3 by $10.2$~pp.
Phase~2b ablation (Table~\ref{tab:bn_ablation}) isolates BatchNorm as the
principal mechanism: removing all eight BatchNorm layers reduces 1D-CNN
accuracy by $12.69$~pp ($83.26\%$), at which point the performance gap
with BM3 ($85.45\%$) is no longer statistically significant ($p = 0.125$).
However, the same statistical dependence that makes BN effective within
distribution makes it brittle under domain shift: on XJTU-SY cross-condition
evaluation, 1D-CNN collapses to a constant-OR predictor
(IR recall $\approx 1\%$, Table~\ref{tab:ood_comparison}).

The \textbf{frequency-grounded paradigm} (BearMamba-3 + $\mathcal{L}_{\text{kin}}$)
replaces statistical normalisation with an explicit, distribution-agnostic
constraint: each training step penalises the distance between learned state
frequencies and analytically known fault harmonics (Section~\ref{subsec:lkin}).
This approach yields no within-distribution advantage at the scale tested
on CWRU (a near-ceiling benchmark), but maintains macro-F1 of $79.73\%$
on XJTU-SY cross-condition --- $38.4$~pp above 1D-CNN (BN) --- because
the physical prior remains valid regardless of source-domain statistics.

The practical implication is a deployment-context decision rather than an
absolute ranking:
\emph{within-condition} monitoring (same machine, constant speed) favours
BN-dependent methods for their stronger implicit regularisation;
\emph{cross-machine or variable-speed} monitoring favours frequency-grounded
methods whose constraint validity is independent of the source domain.
Future architectures might combine both: a hybrid regulariser that activates
BatchNorm within domain and $\mathcal{L}_{\text{kin}}$ at domain boundaries.
This regime-dependent positioning is corroborated by concurrent
physics-guided Mamba work: PG-TMT~\cite{li2026pgtmt} similarly
sidesteps direct comparison against classical 1D-CNN baselines and
positions instead on edge deployment and reliability calibration,
paralleling our framing of physical inductive bias as a
complementary---rather than competitive---axis to within-distribution
accuracy.
% ⚠️ D33 Win2 叙事框架最终版：BN解释机制 + OOD反转；不写"Win2 achieved"

\paragraph{Comparison scope and absolute accuracy context.}
This study focuses on ablating the SSM backbone (BM2 vs BM3) and the
physical regulariser ($\pm\mathcal{L}_{\text{kin}}$), holding all other
design choices constant to isolate the effect of each.
To situate BearMamba-3 in a broader context,
Table~\ref{tab:cwru_baselines} includes comparisons against
SK-SVM~\cite{antoni2006spectral_kurtosis}, 1D-CNN, and Transformer-1D.
At clean and moderate SNR, all deep methods approach ceiling accuracy on
CWRU; at lower SNR, methods with stronger implicit regularisation (1D-CNN
with BatchNorm, BearMamba-2) achieve higher accuracy than BearMamba-3,
consistent with BM3's known complex-state optimisation landscape sensitivity
(Section~\ref{subsec:c1}).
The contribution of BearMamba-3 is not accuracy maximisation on
CWRU within-condition, but the structural and interpretability properties
enabled by Mamba-3's oscillatory states (C1: unique carrier of
$\mathcal{L}_{\text{kin}}$), the variance-stabilising physical inductive bias
(C2), and the OOD robustness demonstrated on XJTU-SY cross-condition
(Section~\ref{subsec:paradigm}).

\paragraph{Noise model scope.}
All noise experiments use Additive White Gaussian Noise (AWGN), which
provides a controlled, reproducible benchmark aligned with prior work on
this dataset~\cite{zhang2020deep}.
Industrial environments often include non-Gaussian noise such as coloured
or impulsive interference; the robustness of $\mathcal{L}_{\text{kin}}$
under such conditions is an open question left to future work.

\paragraph{Reproducibility and statistical power.}
All experiments use 5 independent random seeds; the minimum detectable
one-sided Wilcoxon $p$-value is 0.031 (all-concordant).
At this sample size, a true mean difference of $\sim$2~pp with
$\sigma \approx 1\%$ yields statistical power of only $\approx 50\%$;
the absence of significance for C2 variance reduction and several
C3 SNR points should therefore be interpreted as a power limitation
rather than evidence of no effect.

\paragraph{Window-level split and temporal overlap.}
CWRU experiments use a sliding window with stride $= T/2$ (50\% overlap).
Randomly assigning overlapping windows to train/val may introduce weak
temporal correlation, which could marginally inflate accuracy estimates.
This is consistent with standard CWRU benchmarking practice and is
unlikely to affect the relative comparisons between methods.

\paragraph{Deployment cost versus accuracy trade-off.}
While our BearMamba-3 requires a workstation-class GPU ($\geq 4$\,GB)
for deployment, PG-TMT~\cite{li2026pgtmt} achieves edge-IoT deployment
at sub-1\,MB with sub-7\,ms Jetson p99 latency.
We position this work for industrial PC deployments where cross-condition
accuracy and physical interpretability take precedence over
footprint---a complementary regime to edge-IoT diagnostic instrumentation,
not a competing one.
Future work could investigate quantisation-aware distillation of
BearMamba-3's complex SSM states for edge inference, potentially combined
with EVT-style reliability calibration to bridge the two regimes.
